\subsection*{Appendix A: Mass-Polarization Term in Classical Molecular Hamiltonian}

Section 2 produced the mass-polarization term on the momentum side, via the generating function (2.5). This appendix re-derives it from the velocity side: substitute the same mixed mass-centered coordinates of \S\,2.1 into the laboratory-frame kinetic energy, read off the canonical momenta from the Lagrangian, and Legendre-transform back. Every step is elementary; the payoff is seeing where the mass polarization hides in velocity form --- as a \emph{negative} correction --- and how the Legendre transform flips its sign.

We start from the laboratory-frame kinetic energy
\begin{equation}
T = \frac{1}{2}\sum_j M_j \dot{\mathbf R}_j^2 + \frac{1}{2}\sum_i m_e \dot{\mathbf r}_i^2 , \tag{A.1}
\end{equation}
where $M_j$ is the mass of the $j$-th nucleus, $\mathbf R_j$ its Cartesian coordinate, $m_e$ the electron mass, and $\mathbf r_i$ the coordinate of the $i$-th electron.

\subsubsection*{Mixed Mass-Centered Coordinates}

As in (2.1)--(2.3): the nuclear and total masses and centers of mass are
\begin{equation}
M_N = \sum_j M_j ,\qquad M_{\mathrm{tot}} = M_N + N_{\mathrm{elec}}\,m_e , \tag{A.2}
\end{equation}
\begin{equation}
\mathbf R_{\mathrm{ncm}} = \frac{1}{M_N}\sum_j M_j \mathbf R_j ,\qquad
\mathbf R_{\mathrm{cm}} = \frac{1}{M_{\mathrm{tot}}}\Bigl(\sum_j M_j \mathbf R_j + m_e\sum_i \mathbf r_i\Bigr)
= \mathbf R_{\mathrm{ncm}} + \frac{m_e}{M_{\mathrm{tot}}}\sum_i \mathbf r'_i , \tag{A.3}
\end{equation}
and the internal coordinates are referred to the \emph{nuclear} CoM,
\begin{equation}
\mathbf R'_j = \mathbf R_j - \mathbf R_{\mathrm{ncm}} ,\qquad
\mathbf r'_i = \mathbf r_i - \mathbf R_{\mathrm{ncm}} ,\qquad
\sum_j M_j \mathbf R'_j = 0 . \tag{A.4}
\end{equation}
The constraint ties the nuclear relative vectors only; the $\mathbf r'_i$ are unconstrained. To substitute into $T$ we need the old coordinates in terms of the new. Solving the second equality of the $\mathbf R_{\mathrm{cm}}$ definition for the nuclear CoM,
\begin{equation}
\mathbf R_{\mathrm{ncm}} = \mathbf R_{\mathrm{cm}} - \frac{m_e}{M_{\mathrm{tot}}}\sum_l \mathbf r'_l , \tag{A.5}
\end{equation}
and therefore
\begin{equation}
\mathbf R_j = \mathbf R_{\mathrm{cm}} + \mathbf R'_j - \frac{m_e}{M_{\mathrm{tot}}}\sum_l \mathbf r'_l ,
\qquad
\mathbf r_i = \mathbf R_{\mathrm{cm}} + \mathbf r'_i - \frac{m_e}{M_{\mathrm{tot}}}\sum_l \mathbf r'_l . \tag{A.6}
\end{equation}
Both families ride on the \emph{same} drift --- the velocity of the nuclear CoM,
\begin{equation}
\mathbf D \equiv \dot{\mathbf R}_{\mathrm{ncm}} = \dot{\mathbf R}_{\mathrm{cm}} - \frac{m_e}{M_{\mathrm{tot}}}\sum_l \dot{\mathbf r}'_l ,
\qquad
\dot{\mathbf R}_j = \mathbf D + \dot{\mathbf R}'_j ,
\qquad
\dot{\mathbf r}_i = \mathbf D + \dot{\mathbf r}'_i . \tag{A.7}
\end{equation}

\subsubsection*{Kinetic Energy in Velocity Form}

Substitute the velocities into $T$. The nuclear part, using $\sum_j M_j\dot{\mathbf R}'_j = 0$ (the time derivative of the constraint) on the cross term,
\begin{equation}
T_N = \frac12\sum_j M_j\bigl(\mathbf D + \dot{\mathbf R}'_j\bigr)^2
= \frac12 M_N \mathbf D^2 + \mathbf D\cdot\underbrace{\sum_j M_j\dot{\mathbf R}'_j}_{=\,0} + \frac12\sum_j M_j\dot{\mathbf R}_j^{'2} . \tag{A.8}
\end{equation}
The electronic part has no constraint, so its cross term survives:
\begin{equation}
T_e = \frac12 m_e\sum_i\bigl(\mathbf D + \dot{\mathbf r}'_i\bigr)^2
= \frac12 N_{\mathrm{elec}}\,m_e \mathbf D^2 + m_e\,\mathbf D\cdot\sum_i\dot{\mathbf r}'_i + \frac12 m_e\sum_i\dot{\mathbf r}_i^{'2} . \tag{A.9}
\end{equation}
Adding, with $M_N + N_{\mathrm{elec}}m_e = M_{\mathrm{tot}}$,
\begin{equation}
T = \frac12 M_{\mathrm{tot}}\mathbf D^2 + m_e\,\mathbf D\cdot\sum_i\dot{\mathbf r}'_i + \frac12\sum_j M_j\dot{\mathbf R}_j^{'2} + \frac12 m_e\sum_i\dot{\mathbf r}_i^{'2} . \tag{A.10}
\end{equation}
Now expand $\mathbf D = \dot{\mathbf R}_{\mathrm{cm}} - \frac{m_e}{M_{\mathrm{tot}}}\sum_l\dot{\mathbf r}'_l$ in the first two terms:
\begin{align}
\frac12 M_{\mathrm{tot}}\mathbf D^2
&= \frac12 M_{\mathrm{tot}}\dot{\mathbf R}_{\mathrm{cm}}^2
- m_e\,\dot{\mathbf R}_{\mathrm{cm}}\!\cdot\!\sum_l\dot{\mathbf r}'_l
+ \frac{m_e^2}{2M_{\mathrm{tot}}}\Bigl(\sum_l\dot{\mathbf r}'_l\Bigr)^2 , \tag{A.11a}
\\
m_e\,\mathbf D\cdot\sum_i\dot{\mathbf r}'_i
&= m_e\,\dot{\mathbf R}_{\mathrm{cm}}\!\cdot\!\sum_i\dot{\mathbf r}'_i
- \frac{m_e^2}{M_{\mathrm{tot}}}\Bigl(\sum_l\dot{\mathbf r}'_l\Bigr)^2 . \tag{A.11b}
\end{align}
The two $\dot{\mathbf R}_{\mathrm{cm}}\!\cdot\!\sum\dot{\mathbf r}'$ cross terms cancel \emph{identically} --- the velocity-form image of the momentum-side cancellation in \S\,2.3 --- and the quadratic pieces combine to a single negative remainder:
\begin{equation}
\boxed{\;
T = \frac12 M_{\mathrm{tot}}\dot{\mathbf R}_{\mathrm{cm}}^2
+ \frac12\sum_j M_j\dot{\mathbf R}_j^{'2}
+ \frac12 m_e\sum_i\dot{\mathbf r}_i^{'2}
- \frac{m_e^2}{2M_{\mathrm{tot}}}\Bigl(\sum_i\dot{\mathbf r}'_i\Bigr)^2 .
\;} \tag{A.12}
\end{equation}
Two features are already visible at the velocity level: the CoM decouples with the \emph{correct} total mass $M_{\mathrm{tot}}$, and the electron sector carries a \emph{negative} collective correction. That negative term is the mass polarization in disguise --- the Legendre transform below flips its sign and trades $M_{\mathrm{tot}}$ for $M_N$.

\subsubsection*{Canonical Momenta}

With $L = T - V$ and $V$ velocity-independent, differentiate the boxed $T$.

\noindent\emph{CoM:} only the first term carries $\dot{\mathbf R}_{\mathrm{cm}}$, so
\begin{equation}
\boxed{\;\mathbf P_{\mathrm{cm}} = \frac{\partial L}{\partial\dot{\mathbf R}_{\mathrm{cm}}} = M_{\mathrm{tot}}\dot{\mathbf R}_{\mathrm{cm}} .\;} \tag{A.13}
\end{equation}
Sanity check: $\sum_j M_j\dot{\mathbf R}_j + m_e\sum_i\dot{\mathbf r}_i = M_{\mathrm{tot}}\mathbf D + m_e\sum_i\dot{\mathbf r}'_i = M_{\mathrm{tot}}\dot{\mathbf R}_{\mathrm{cm}}$, so $\mathbf P_{\mathrm{cm}}$ is the total momentum, as it must be.

\noindent\emph{Electrons:} the last two terms of the boxed $T$ carry $\dot{\mathbf r}'_i$,
\begin{equation}
\boxed{\;\mathbf p'_i = \frac{\partial L}{\partial\dot{\mathbf r}'_i}
= m_e\dot{\mathbf r}'_i - \frac{m_e^2}{M_{\mathrm{tot}}}\sum_l\dot{\mathbf r}'_l ,\;} \tag{A.14}
\end{equation}
which is precisely the momentum--velocity relation (3.5). Summing over $i$,
\begin{equation}
\sum_i\mathbf p'_i = m_e\Bigl(1 - \frac{N_{\mathrm{elec}}m_e}{M_{\mathrm{tot}}}\Bigr)\sum_l\dot{\mathbf r}'_l
= \frac{m_e M_N}{M_{\mathrm{tot}}}\sum_l\dot{\mathbf r}'_l , \tag{A.15}
\end{equation}
so $\frac{m_e^2}{M_{\mathrm{tot}}}\sum_l\dot{\mathbf r}'_l = \frac{m_e}{M_N}\sum_l\mathbf p'_l$, and the relation inverts to
\begin{equation}
m_e\dot{\mathbf r}'_i = \mathbf p'_i + \frac{m_e}{M_N}\sum_l\mathbf p'_l
\qquad\Longleftrightarrow\qquad
\dot{\mathbf r}'_i = \frac{\mathbf p'_i}{m_e} + \frac{1}{M_N}\sum_l\mathbf p'_l , \tag{A.16}
\end{equation}
the Hamilton equation (3.4).

\noindent\emph{Nuclei:} on the constraint surface $\sum_j M_j\dot{\mathbf R}'_j = 0$ and in the dual gauge $\sum_j\mathbf P'_j = 0$ of \S\,2.2,
\begin{equation}
\boxed{\;\mathbf P'_j = M_j\dot{\mathbf R}'_j .\;} \tag{A.17}
\end{equation}
Consistency check against (2.8): the lab momentum is $M_j\dot{\mathbf R}_j = M_j\mathbf D + M_j\dot{\mathbf R}'_j$, and
\begin{equation}
\mathbf D = \dot{\mathbf R}_{\mathrm{cm}} - \frac{m_e}{M_{\mathrm{tot}}}\sum_l\dot{\mathbf r}'_l
= \frac{\mathbf P_{\mathrm{cm}}}{M_{\mathrm{tot}}} - \frac{1}{M_N}\sum_l\mathbf p'_l , \tag{A.18}
\end{equation}
which is exactly the common vector $\mathbf b$ of (2.9); hence $\mathbf P_j = \frac{M_j}{M_{\mathrm{tot}}}\mathbf P_{\mathrm{cm}} + \mathbf P'_j - \frac{M_j}{M_N}\sum_l\mathbf p'_l$, reproducing (2.8).

\subsubsection*{Kinetic Energy in Momentum Form}

The CoM and nuclear terms are immediate: $\frac12 M_{\mathrm{tot}}\dot{\mathbf R}_{\mathrm{cm}}^2 = \mathbf P_{\mathrm{cm}}^2/2M_{\mathrm{tot}}$ and $\frac12 M_j\dot{\mathbf R}_j^{'2} = \mathbf P_j^{'2}/2M_j$. For the electron sector, abbreviate $\bar{\mathbf P}\equiv\sum_l\mathbf p'_l$ and substitute $\dot{\mathbf r}'_i = \mathbf p'_i/m_e + \bar{\mathbf P}/M_N$:
\begin{equation}
\frac12 m_e\sum_i\dot{\mathbf r}_i^{'2}
= \sum_i\frac{1}{2m_e}\Bigl(\mathbf p'_i + \frac{m_e}{M_N}\bar{\mathbf P}\Bigr)^2
= \sum_i\frac{\mathbf p_i^{'2}}{2m_e} + \frac{\bar{\mathbf P}^2}{M_N} + \frac{N_{\mathrm{elec}}m_e}{2M_N^2}\bar{\mathbf P}^2 , \tag{A.19}
\end{equation}
and, using $\sum_l\dot{\mathbf r}'_l = \frac{M_{\mathrm{tot}}}{m_e M_N}\bar{\mathbf P}$,
\begin{equation}
-\frac{m_e^2}{2M_{\mathrm{tot}}}\Bigl(\sum_l\dot{\mathbf r}'_l\Bigr)^2
= -\frac{m_e^2}{2M_{\mathrm{tot}}}\cdot\frac{M_{\mathrm{tot}}^2}{m_e^2 M_N^2}\,\bar{\mathbf P}^2
= -\frac{M_{\mathrm{tot}}}{2M_N^2}\,\bar{\mathbf P}^2 . \tag{A.20}
\end{equation}
Collecting the three $\bar{\mathbf P}^2$ coefficients,
\begin{equation}
\frac{1}{M_N} + \frac{N_{\mathrm{elec}}m_e - M_{\mathrm{tot}}}{2M_N^2}
= \frac{1}{M_N} - \frac{M_N}{2M_N^2}
= \frac{1}{M_N} - \frac{1}{2M_N}
= +\frac{1}{2M_N} , \tag{A.21}
\end{equation}
using $M_{\mathrm{tot}} - N_{\mathrm{elec}}m_e = M_N$. The negative velocity-form correction has become the \emph{positive} mass polarization:
\begin{equation}
\boxed{\;
T = \frac{\mathbf P_{\mathrm{cm}}^2}{2M_{\mathrm{tot}}}
+ \sum_j\frac{\mathbf P_j^{'2}}{2M_j}
+ \sum_i\frac{\mathbf p_i^{'2}}{2m_e}
+ \frac{1}{2M_N}\Bigl(\sum_i\mathbf p'_i\Bigr)^2 .
\;} \tag{A.22}
\end{equation}

\subsubsection*{Hamiltonian}

The potential is translationally invariant and both internal families share the reference point $\mathbf R_{\mathrm{ncm}}$, so every interparticle separation is a plain coordinate difference and $V = V(\{\mathbf R'_j\},\{\mathbf r'_i\})$. Hence
\begin{equation}
\boxed{\;
H = \frac{\mathbf P_{\mathrm{cm}}^2}{2M_{\mathrm{tot}}}
+ \sum_j\frac{\mathbf P_j^{'2}}{2M_j}
+ \sum_i\frac{\mathbf p_i^{'2}}{2m_e}
+ \frac{1}{2M_N}\Bigl(\sum_i\mathbf p'_i\Bigr)^2
+ V(\{\mathbf R'_j\},\{\mathbf r'_i\}) ,
\;} \tag{A.23}
\end{equation}
identical to (2.12) --- exact for \emph{any} $\mathbf P_{\mathrm{cm}}$, with no rest-frame choice, no CoM--internal cross-coupling, and the correct total mass in the free term.

\subsubsection*{Important Interpretation}

\noindent\emph{The canonical electron momentum is not $m_e\dot{\mathbf r}'_i$.} It is depleted by the collective recoil,
\begin{equation}
\mathbf p'_i = m_e\dot{\mathbf r}'_i - \frac{m_e^2}{M_{\mathrm{tot}}}\sum_l\dot{\mathbf r}'_l , \tag{A.24}
\end{equation}
because moving one electron drags the total CoM, and the coordinate $\mathbf r'_i$ is measured from the (recoiling) nuclear CoM.

\noindent\emph{Why the sign flips.} The electron block of the velocity-form kinetic energy is the quadratic form $\frac12\dot{\mathbf r}^{\prime\mathsf T}\mathbf A\,\dot{\mathbf r}'$ with the non-diagonal mass matrix
\begin{equation}
\mathbf A = m_e\Bigl(\mathbb{1} - \frac{m_e}{M_{\mathrm{tot}}}\mathbf J\Bigr),
\qquad \mathbf J_{il} = 1 \ \text{(all-ones block)}, \tag{A.25}
\end{equation}
and the Legendre transform replaces $\mathbf A$ by its \emph{inverse}:
\begin{equation}
\mathbf A^{-1} = \frac{1}{m_e}\Bigl(\mathbb{1} + \frac{m_e}{M_N}\mathbf J\Bigr)
\qquad\Longrightarrow\qquad
\frac12\,\mathbf p^{\prime\mathsf T}\mathbf A^{-1}\mathbf p'
= \sum_i\frac{\mathbf p_i^{'2}}{2m_e} + \frac{1}{2M_N}\Bigl(\sum_i\mathbf p'_i\Bigr)^2 , \tag{A.26}
\end{equation}
(one verifies $\mathbf A\mathbf A^{-1}=\mathbb{1}$ using $\mathbf J^2 = N_{\mathrm{elec}}\mathbf J$ and $M_{\mathrm{tot}} - N_{\mathrm{elec}}m_e = M_N$). The inversion turns the negative $-m_e^2/2M_{\mathrm{tot}}$ collective coupling into the positive $+1/2M_N$ mass polarization: same physics, opposite side of the Legendre transform.

\noindent\emph{Bookkeeping.} The accompanying constraints are
\begin{equation}
\boxed{\;
\sum_j M_j \mathbf R'_j = 0 ,
\qquad
\sum_j \mathbf P'_j = 0 ,
\;} \tag{A.27}
\end{equation}
purely nuclear on both sides. Everything above agrees with the generating-function route of Section 2: the momentum relations reproduce (2.8), and the Hamiltonian reproduces (2.12).
